% science_template.tex
% See accompanying readme.txt for copyright statement, change log etc.

% Any modification of this template, including writing a paper using it,
% MUST rename the file i.e. use a different file name.

%%%%%%%%%%%%%%%% START OF PREAMBLE %%%%%%%%%%%%%%%

% Basic setup. Authors shouldn't need to adjust these commands.
% It's annoying, but please do NOT strip these into a separate file.
% They need to be included in this .tex for our production software to work.

% Use the basic LaTeX article class, 12pt text
\documentclass[12pt]{article}

% Science uses Times font. If you don't have this installed (most LaTeX installations will be
% fine) or prefer the old Computer Modern fonts, comment out the following line
\usepackage{newtxtext,newtxmath}
% Depending on your LaTeX fonts installation, you might get better results with one or both of these:
%\usepackage{mathptmx}
%\usepackage{txfonts}

% Allow external graphics files
\usepackage{graphicx}

%%% Tompkins IMPORTS
% Mathematics and Times-style fonts
\usepackage{amsmath}

% Use US letter sized paper with 1 inch margins
\usepackage[letterpaper,margin=1in]{geometry}

%%% Tompkins imports
\usepackage{tikz}
\usepackage{newfloat}
\usetikzlibrary{decorations.pathreplacing}

\setlength{\oddsidemargin}{0in}
\setlength{\evensidemargin}{0in}
\setlength{\textwidth}{6.5in}
\setlength{\topmargin}{-0.5in}
\setlength{\textheight}{9in}
\setlength{\parindent}{0pt}
\setlength{\parskip}{0.9em}
\setlength{\emergencystretch}{3em}

\newcommand{\photoplaceholder}[1]{%
    \begin{center}
    \fbox{\begin{minipage}[c][2.2in][c]{0.78\textwidth}
    \centering\textbf{PHOTO PLACEHOLDER}\\[8pt]#1
    \end{minipage}}
    \end{center}}

% Drawings receive their own “Example” counter.
\DeclareFloatingEnvironment[
    name=Example,
    listname={List of Examples}
]{example}

% Double line spacing, including in captions
\linespread{1.5} % For some reason double spacing is 1.5, not 2.0!

% One space after each sentence
\frenchspacing

% Abstract formatting and spacing - no heading
\renewenvironment{abstract}
{\quotation}
{\endquotation}

% No date in the title section
\date{}

% Reference section heading
\renewcommand\refname{References and Notes}

% Figure and Table labels in bold
\makeatletter
\renewcommand{\fnum@figure}{\textbf{Figure \thefigure}}
\renewcommand{\fnum@table}{\textbf{Table \thetable}}
\makeatother



% Provides the \url command, and fixes a crash if URLs or DOIs contain underscores
\usepackage{url}

%%%%%%%%%%%% CUSTOM COMMANDS AND PACKAGES %%%%%%%%%%%%

% Authors can define simple custom commands e.g. as shortcuts to save on typing
% Use \newcommand (not \def) to avoid overwriting existing commands.
% Keep them as simple as possible and note the warning in the text below.

% Example:
\newcommand{\pcc}{\,cm$^{-3}$}	% per cm-cubed

% If your BibTeX uses shortcuts for journal names, define them here e.g.
% \newcommand{\apj}{Astrophys.~J.}

% Please DO NOT import additional external packages or .sty files.
% Those are unlikely to work with our conversion software and will cause problems later.
% Don't add any more \usepackage{} commands.


%%%%%%%%%%%%%%%% TITLE AND AUTHORS %%%%%%%%%%%%%%%%

% Title of the paper.
% Keep it short and understandable by any reader of Science.
% Avoid acronyms or jargon. Use sentence case.
\def\scititle{
    \textbf{A SIMPLE PROOF OF P v NP}
}
% Store the title in a variable for reuse in the supplement (otherwise \maketitle deletes it)
\title{\bfseries \boldmath \scititle}

% Author and institution list.
% Institution numbers etc. should be hard-coded, do *not* use the \footnote command.
\author{
% You can write out first names or use initials - either way is acceptable, but be consistent
    Christopher Robinson Tompkins III, Independent, Richmond, VA\and\
% Do not include street addresses, postal codes etc.
%
% Identify at least one corresponding author, with contact email address
    Corresponding author. Email: chtompki@apache.org,\\ chtompki@vt.edu, tompkinscr@vcu.edu
% Joint contributions can be indicated like this
}

%%%%%%%%%%%%%%%%% END OF PREAMBLE %%%%%%%%%%%%%%%%


%%%%%%%%%%%%%%%% START OF MAIN TEXT %%%%%%%%%%%%%%%
\begin{document}

% Insert the title and author list
    \maketitle

% Abstract, in bold
% There are strict length limits, and not all formats have abstracts.
% Consult the journal instructions to authors for details.
% Do not cite any references in the abstract.
    \begin{abstract} \bfseries \boldmath
        A proof of the PvNP problem such that a grade schooler should be able to understand. I am
        writing for my daughters Margo (10) and Nora (12).
    \end{abstract}


    \textbf{(0)}
    \begin{center}
        \emph{Dedicated to the memory of those killed at Virginia Tech,\\
        and in solidarity with those who were injured, their families,\\
        and the Virginia Tech community.}
        \qquad\\
        \emph{I would like to acknowledge the following people that helped me along the way: mom, dad, Margo, Nora,
            Griff Elder (and the 2007 VT Mathematics department and all my friends from there), George McNulty (and the 2009 USC Mathematics department and all my friends from there),
            Andy Lewis (and the VCU 2026 Mathematics Department and all my friends from there), and those at NASA whom I've worked with.}
    \end{center}

    \newpage
    \textbf{(1)}
    \section*{DEFINITIONS}

    Before going into the proof we need to clarify what the complexity classes \textbf{P} and \textbf{NP} are.

    \section*{REFERENCES}

    \begin{enumerate}
        \item Stephen Cook, ``The complexity of theorem-proving procedures,''
        \emph{Conference Record of Third Annual ACM Symposium on Theory of Computing},
        pages 151-158, 1971.
        \item Stephen Cook, ``The \bf{P} versus \bf{NP} Problem,''
        \emph{Official Problem Description for the Millennium Prize Problems},
        Clay Mathematics Institute, 2000,
        \url{https://www.claymath.org/wp-content/uploads/2022/06/pvsnp.pdf}
    \end{enumerate}

\end{document}